\documentclass{article}
\usepackage{float}
\usepackage{textcomp}
\usepackage{graphicx}
\usepackage{booktabs}
\usepackage{color}
\usepackage{verbatim}
\usepackage{listings}
\usepackage{underscore}
\usepackage{amsmath}
\usepackage{amssymb}
\usepackage{listings}
%\usepackage{showkeys}
%\usepackage{float} % here for H placement parameter
\usepackage{flafter} 
\setcounter{secnumdepth}{5}
\usepackage[bookmarks=true]{hyperref}
\usepackage[dvipsnames]{xcolor}
\author{Mirko Mantovani (893784), Matteo Marziali (893904)} 
\date{\today}
\title{Politecnico di Milano
	\\A.A. 2017\@-\@2018
	\\Software Engineering II project \\ \textbf{Travlendar+}
	\\
	\textbf{A}cceptance \textbf{T}est 
	\\ \textbf{D}ocument 
	\\
	\textbf{V1}} 
	\hypersetup{pdftitle={Acceptance Test Document},    % title
	pdfauthor={Mirko Mantovani, Matteo Marziali},                     % author
	pdfsubject={ATD},                        % subject of the document
	colorlinks=true,       % false: boxed links; true: colored links
	linkcolor=blue,       % color of internal links
	citecolor=blue,       % color of links to bibliography
	filecolor=black,        % color of file links
	urlcolor=purple,        % color of external links
   }
\begin{document}
\maketitle
\begin{center}
	\includegraphics[width=7cm]{polimi-logo}
\end{center}
\clearpage
{\hypersetup{hidelinks}\tableofcontents}


\clearpage
\section{Analyzed project}
\subsection{Project info}
The project analysed is \textbf{EventMI} developed by: Guglielmo Manneschi and Andrea Mecchia.
\\The link to the private repository containing code and documentation is: \\
\href{url}{https://github.com/nondecidibile/ManneschiMecchia}

\subsection{Documents considered}
The documents we took into consideration when analysing and testing the application are:
\begin{itemize}
\item \textbf{ITD:} from which we found the implemented requirements/functionalities.
\item \textbf{RASD1.1:} from which we found all the requirements and checked the coherence with the ones that were actually implemented.

\end{itemize}


\clearpage
\section{Installation Setup}
Since the last version of the application was already remotely deployed and the group provided users with public access to it at the address page: \href{url}{http://eventmi.date}, we avoided the installation on our machines and preferred to concentrate the little time we had on performing a thorough testing of the requirements they implemented.\\ The installation manual seems clear though, even if a installation guide based on a more widespread operating system, such as Windows, would have been better.

\clearpage
\section{Acceptance Test Cases}
\subsection{Login/Signup}

\begin{flushleft}

\begin{table}[htp]

\begin{tabular}{p{7cm}|p{7cm}}
 Implemented requirement&Notes\\
\hline
\hline
R1: Registration&\textcolor{green}{Accomplished}\\
\hline
R2: Login&\textcolor{green}{Accomplished}\\
\hline

\end{tabular}

\caption{Req table } 
\label{tab:Reqtable}

\end{table}

\end{flushleft}



\subsection{Event CRUD operations}

\begin{flushleft}
	
	\begin{table}[htp]
		
		\begin{tabular}{p{7cm}|p{7cm}}
			Implemented requirement&Notes\\
			\hline
			\hline

R3: The insertion of new event in the calendar. & \textcolor{green}{Accomplished}. See \hyperlink{meetingRemark}{Additional notes.}\\
\hline
R4: Modification of an event. & \textcolor{green}{Accomplished}\\
\hline
R5: Deletion of an event. & \textcolor{green}{Accomplished}\\
\hline
R7: Warning creation. & \textcolor{BurntOrange}{Partially accomplished} Everything works fine excepts for a minor bug. We made a test case consisting in inserting a meeting from 22.00p.m. to 23.00p.m. in Busto Arsizio and a meeting from 2.00a.m. to 3.00a.m. in Puglia on the following day, obviously the insertion should be prevented due to the impossibility to cover the distance between the meetings in just 3 hours, but this doesn't happen, the second meeting is accepted by the system.  \\
\hline

\end{tabular}

\caption{Event table } 
\label{tab:Eventtable}

\end{table}

\end{flushleft}



\subsection{Routes}

\begin{flushleft}
	
	\begin{table}[htp]
		
		\begin{tabular}{p{7cm}|p{7cm}}
			Implemented requirement&Notes\\
			\hline
			\hline
R8: The identification of the best mobility option for a transfer, according to user preferences and other external factors.& \textcolor{green}{  Accomplished} \\
\hline

\end{tabular}

\caption{Mobility table } 
\label{tab:Mobilitytable}

\end{table}

\end{flushleft}



\subsection{Breaks}

\begin{flushleft}
	
	\begin{table}[htp]
		
		\begin{tabular}{p{7cm}|p{7cm}}
			Implemented requirement&Notes\\
			\hline
			\hline
Definition of flexible daily breaks&\textcolor{green}{Accomplished}\\
\hline
Modification of flexible daily breaks&\\
\hline
Usage of flexible daily breaks&\\
\hline

\end{tabular}

\caption{Break table } 
\label{tab:Breaktable}

\end{table}

\end{flushleft}

\clearpage

\subsection{Travel mean preferences}

\begin{flushleft}
	
	\begin{table}[htp]
		
		\begin{tabular}{p{7cm}|p{7cm}}
			Implemented requirement&Notes\\
			\hline
			\hline
R14: Some settings regarding the navigation profile, such as choosing the preferred means of transport&\textcolor{BurntOrange}{Partially accomplished} The unique preference the user can manage regards the travel means and in particular a prevalence order of them. This leads to a low level of customization for the users. Moreover the prevalence order prevent a user avoiding to adopt a specific travel mean, indeed it is not excluded but at most it is placed below other transports.   \\
\hline

\end{tabular}

\caption{Nav table } 
\label{tab:Navtable}

\end{table}

\end{flushleft}


\clearpage
\subsection{Use cases}
In this section we check that all use cases presented in the RASD matches the implementation.

\begin{flushleft}

\begin{table}[htp]

\begin{tabular}{p{3cm}|p{5cm}|p{5cm}}
\bf\large Use Case&\bf\large Exit Condition&\bf\large Exceptions\\
\hline
\hline
\textbf{Registration}&\large\bf\textcolor{green}{Correct}&\large\bf\textcolor{green}{Correct}\\
\hline
\textbf{Login}&\large\bf\textcolor{green}{Correct}&\large\bf\textcolor{green}{Correct}\\
\hline
\textbf{Add event to calendar}&\large\bf\textcolor{green}{Correct}&\large\bf\textcolor{green}{Correct}\\
\hline
\textbf{Edit an event in the calendar}&\large\bf\textcolor{green}{Correct}&\large\bf\textcolor{green}{Correct}\\
\hline
\textbf{Removing an event from the calendar}&\large\bf\textcolor{green}{Correct}&\large\bf\textcolor{green}{Correct}\\
\hline
\textbf{Recompute transfer for an event}&\large\bf\textcolor{red}{Incorrect}: It displays an error: "recompute not allowed over insecure network", yet the application is not present in a secure domain (https)&\large\bf\textcolor{red}{Incorrect}: It is stated no exceptions presents where instead an exception is always presented\\
\hline
\textbf{Transfer notification}&\large\bf Not implemented&\large\bf Not implemented\\
\hline
\textbf{Define a new break}&\large\bf\textcolor{green}{Correct}&\large\bf\textcolor{green}{Correct}\\
\hline
\textbf{Modify a break}&\large\bf\textcolor{green}{Correct}&\large\bf\textcolor{green}{Correct}\\
\hline
\textbf{Delete a break}&\large\bf\textcolor{green}{Correct}&\large\bf\textcolor{green}{Correct}\\
\hline
\textbf{Add a preferred mean of transport}&\large\bf\textcolor{green}{Correct}&\large\bf\textcolor{green}{Correct}\\
\hline
\textbf{Select a navigation profile
}&\large\bf Not implemented&\large\bf Not implemented\\
\hline




\end{tabular}

\caption{Use cases tests} 
\label{tab:usecases}

\end{table}

\end{flushleft}



\clearpage
\section{Additional notes}
\begin{itemize}
	\item \textbf{Remark about meeting insertion}\\
	Even if the meeting insertion requirement is practically satisfied and the related goal accomplished, i.e. the operation can be done properly, we noticed that the application allows a conceptually wrong usage of this function, indeed it is possible to schedule meeting also backward in the past. However, considering that meetings in the past are obviously unfeasible, this kind of possibility should be prevented to the user, in fact it may cause slowdowns or efficiency issues, for example in case of a huge number of appointments are registered in the calendar. 
	\hypertarget{meetingRemark}\\
	
\end{itemize}



\clearpage
\section{Appendix}

\subsection{Effort spent}
\begin{itemize}

\item Matteo Marziali working hours:  $\approxeq$ 5 hours

\item Mirko Mantovani working hours:  $\approxeq$ 5 hours



\end{itemize}



\end{document}
